\subsection{ECAのSolver構造に関する理論的考察}

ここまでに示した定理および命題によって、
Solver集合 $\mathcal{S}$、評価空間 $X$、
および一般解 $\mathcal{G}$ の間に成立する数学的関係の
一部は形式的に記述することができる。

一方で、これらの構造がなぜECAにおいて自然に現れるのか、
すなわち「なぜECAのSolver構造からこのような一般解が
導かれるのか」という問いについては、
必ずしも同じ形式の証明として記述することが適切とは限らない。

この問題を考察するためには、ECAの構成に至った動機、
特に数学的命題がどのような公理体系および評価体系の上に
成立するのかという問題を考える必要がある。

\subsubsection{「$1+1=2$」という問題から見た評価体系}

「$1+1=2$」という命題は、一見すると自明な算術的事実に
見える。しかし、数学的に厳密に考えれば、この等式の意味は
数の構成、加法の定義、およびそれらを支える公理体系から
切り離して理解することはできない。

したがって、ある命題
\[
P
\]
を考える場合、その真偽あるいは成立性だけでなく、

\[
(\mathcal{A},P)
\]

という形で、命題 $P$ と、それを評価する公理体系
$\mathcal{A}$ の組を考える必要がある。

ECAにおけるSolverは、このような「命題を評価するための
体系」の差異を一般化して取り扱う対象として解釈できる。

\subsubsection{連続体仮説から見た公理体系への依存性}

連続体仮説を考える場合には、さらに明確に
「ある命題の数学的位置づけは、それを評価する公理体系との
関係から決定される」という問題が現れる。

この観点から見ると、数学的対象を単独で扱うのではなく、

\[
\text{対象}
\quad+\quad
\text{評価体系}
\]

として扱う必要性が生じる。

ECAにおけるSolver集合 $\mathcal{S}$ は、
このような評価体系の差異を持つ対象を一つの集合として
取り扱うための構造として考えることができる。

\subsubsection{Solver集合と一般解}

以上の考察から、ECAにおける一般解

\[
\mathcal{G}:\mathcal{S}\longrightarrow X
\]

は、単に個々のSolverに対する解を列挙する写像としてではなく、
異なる評価体系において得られる解を共通の評価空間 $X$ 上で
取り扱うための写像として解釈することができる。

したがって、

\[
\mathcal{S}
\overset{\mathcal{G}}{\longrightarrow}
X
\]

という構造は、異なる公理的・評価的条件のもとにある
数学的対象を、共通の数学的構造へ写像するという
ECAの基本的発想を反映している。

\subsubsection{本節の位置づけ}

ここで述べた議論は、前節までに示した定理の形式的証明を
構成するものではない。

むしろ、

\[
\text{公理体系への依存性}
\longrightarrow
\text{Solverという評価体系の一般化}
\longrightarrow
\text{評価空間 $X$}
\longrightarrow
\text{一般解 $\mathcal{G}$}
\]

というECAの理論的発想を明確化するための論考である。

したがって、本節で述べた内容をもって、
Solver構造から一般解の存在や一意性を直接証明したものとは
みなさない。

形式的な証明については、前節までに確立された定義、
命題および定理のみに基づいて行うものとする。
